\documentclass{beamer}
\usetheme{Madrid}
\usecolortheme{default}

\title[Toronto Bikeshare Demand Forecasting]
{Group 21}

\subtitle{Toronto Bikeshare Demand Forecasting}

\author[Group 21] 
{Abdul Mohammed\inst{1} \and Axel Tang\inst{2} \and Esam Uddin\inst{3}}

\institute[ ] % (optional)
{
  \inst{1}%
  The Edward S. Rogers Sr. Department of Electrical and Computer Engineering
}

\date[IntroML] % (optional)
{Introduction to Machine Learning -- Winter 2026}

\logo{
\includegraphics[height=0.8cm]{logo_uoft}}

\definecolor{uoftblue}{RGB}{6,41,88}
\setbeamercolor{titlelike}{bg=uoftblue}
\setbeamerfont{title}{series=\bfseries}

\begin{document}

\frame{\titlepage}


\begin{frame}
\frametitle{Project Overview}

\begin{itemize}
  \setlength\itemsep{0.6em}
  \item \textbf{Goal:} Forecast hour based bikeshare demand across Toronto.
  \item \textbf{Data:} Ridership (Jan 2025 -- Feb 2026) + station metadata + hourly weather; chronological train/val/test split.
  \item \textbf{Design choice:} K-means ($K=8$) on station coordinates cuts zero inflation from 77.7\% to 37.3\%.
  \item \textbf{Models:} Baselines with Linear Regression, improve with Random Forest, XGBoost and PyTorch MLP.
\end{itemize}

\end{frame}

\begin{frame}
\frametitle{Data Pipeline, EDA, and Feature Engineering}

{\footnotesize
\begin{itemize}
  \item \textbf{Ingestion:} 12 CSVs; trips outside 1\,min to 24\,hr dropped; aggregated to station-hour pickups with zeros filled.
  \item \textbf{Zero Inflation with Clustering:} 77.7\% station-hours are zero; K-means ($K=8$) on lat/lon groups 999 stations into 8 clusters, cutting zeros to \textbf{37.3\%}.
  \item \textbf{Features \& Split:} 12 features incl.\ \texttt{lag\_24h}; 2025 train $|$ Jan 2026 val $|$ Feb 2026 test.
\end{itemize}
}

\vfill

\centering
\includegraphics[width=0.49\linewidth,height=0.43\textheight,keepaspectratio]{results/eda_cluster_map.png}\hspace{0.01\linewidth}\includegraphics[width=0.49\linewidth,height=0.43\textheight,keepaspectratio]{results/eda_heatmap.png}

\vspace{0.5cm}



\end{frame}





\begin{frame}
\frametitle{Baseline: Linear Regression}

\begin{columns}[T,onlytextwidth]

\column{0.52\textwidth}
{\footnotesize
\begin{tabular}{lcc}
\textbf{Metric} & \textbf{Train} & \textbf{Test} \\
\hline
MAE   & 30.71 & 11.18 \\
RMSE  & 93.52 & 33.01 \\
$R^2$ & 0.89  & 0.71  \\
\end{tabular}
}

\vspace{0.3cm}

{\footnotesize
\begin{itemize}
  \item \textbf{lag\_24h:} Standard forecasting practice; 24h chosen over 1h to avoid handing the model the answer.
  \item \textbf{Limits:} LR fits one slope per feature, cannot separate summer vs.\ winter rush hours; tree models and NNs close this gap.
\end{itemize}
}

\column{0.45\textwidth}
\centering
\includegraphics[width=\linewidth,height=0.65\textheight,keepaspectratio]{results/baseline_pred_vs_actual.png}
{\tiny Linear Regression -- predicted vs actual (test set).}

\end{columns}
\end{frame}


\begin{frame}
\frametitle{Results: Model Comparison}

\centering
{\scriptsize
\setlength{\tabcolsep}{4pt}
\begin{tabular}{lccc}
\textbf{Model} & \textbf{MAE} & \textbf{RMSE} & \textbf{$R^2$} \\
\hline
Linear Regression     & 11.18 & 33.01 & 0.708 \\
Random Forest         & 7.48  & 23.44 & 0.853 \\
XGBoost               & \textbf{7.15}  & \textbf{22.03} & \textbf{0.870} \\
PyTorch MLP           & 7.61  & 23.52 & 0.852 \\
\end{tabular}
}

\vspace{0.1cm}

{\scriptsize
\begin{itemize}
  \item Tuned Random Forest, XGBoost, and PyTorch MLP on the validation set.
  \item XGBoost performed best, ensemble tree methods are the best fit for this tabular forecasting task.
\end{itemize}
}

\vspace{0.1cm}

\centering
\begin{tabular}{cc}
  \includegraphics[width=0.46\linewidth,height=0.38\textheight,keepaspectratio]{results/best_tuned_tree_pred_vs_actual.png} &
  \includegraphics[width=0.46\linewidth,height=0.38\textheight,keepaspectratio]{results/pytorch_mlp_loss_curve.png} \\[-2pt]
  {\tiny XGBoost (tuned) -- predicted vs actual.} &
  {\tiny PyTorch MLP -- predicted vs actual.}
\end{tabular}

\end{frame}




\end{document}